\documentclass[12pt]{article}

\usepackage{sbc-template}
\usepackage{graphicx,url}
\usepackage[utf8]{inputenc}
\usepackage[brazil]{babel}

\sloppy

\title{Biblioteca para Manipulação de Grafos Simples e Direcionados:\\Relatório da Etapa 1}

\author{Luan  Inaibes\inst{1}, Otávio Santos\inst{1}, Arthur Luis\inst{1}}

\address{Curso de Ciência de Dados e IA -- Pontifícia Universidade Católica de Minas Gerais \\
  Belo Horizonte -- Minas Gerais -- Brasil
}

\begin{document}

\maketitle

\begin{abstract}

\end{abstract}

\begin{resumo}

\end{resumo}

\section{Arquitetura e API}

A biblioteca é estruturada em torno da classe abstrata
\texttt{AbstractGraph}, que define a API comum, mantém os atributos
compartilhados (como os pesos dos vértices), fornece métodos auxiliares
(por exemplo, a validação de índices de vértices) e declara os métodos
abstratos implementados pelas classes concretas. Algoritmos
independentes de representação, como a verificação de conectividade e a
exportação para o Gephi, são implementados uma única vez na classe
abstrata.

Duas classes concretas estendem \texttt{AbstractGraph}:

\begin{itemize}
  \item \texttt{AdjacencyMatrixGraph}: representa as arestas em uma matriz
  $n \times n$, em que cada posição armazena o peso da aresta
  correspondente, favorecendo consultas de existência e de peso em tempo
  constante;
  \item \texttt{AdjacencyListGraph}: representa as arestas por listas de
  adjacência (dicionários destino~$\rightarrow$~peso), com conjuntos de
  predecessores auxiliares, favorecendo grafos esparsos.
\end{itemize}

Ambas recebem no construtor o número de vértices do grafo e respeitam o
mesmo contrato comportamental. A API inclui, entre outros métodos:
contagem de vértices e arestas; \texttt{addEdge}, \texttt{removeEdge} e
\texttt{hasEdge}; consultas de sucessor e predecessor; predicados
estruturais (\texttt{isDivergent}, \texttt{isConvergent},
\texttt{isIncident}); cálculo de grau de entrada e de saída; definição e
consulta de pesos de vértices e de arestas.

\section{Restrições e Tratamento de Erros}

A operação \texttt{addEdge} é idempotente (chamadas repetidas com o mesmo par de
vértices não criam arestas duplicadas) e proíbe laços. São lançadas
exceções específicas para índices de vértices inválidos e para operações
inconsistentes.

\section{Demonstração}

Acompanha a biblioteca um programa de demonstração, executável em
terminal, que percorre as principais funcionalidades da API construindo o
mesmo grafo de exemplo nas duas representações e comparando os resultados
lado a lado: criação dos grafos, inserção e remoção de arestas, consulta
de sucessores e predecessores, cálculo de graus, definição e consulta de
pesos de vértices e de arestas, verificações de grafo vazio, completo e
conectado, e exportação para o Gephi.
\section{Divisão do Trabalho}

A Tabela~\ref{tab:responsabilidades} apresenta a divisão de
responsabilidades entre os integrantes do grupo nesta etapa.

\begin{table}[ht]
\centering
\caption{Divisão de responsabilidades do trabalho.}
\label{tab:responsabilidades}
\begin{tabular}{|l|p{8cm}|}
\hline
\textbf{Nome} & \textbf{Responsabilidades do trabalho} \\
\hline
Luan Inaibes & Implementação da classe abstrata \texttt{AbstractGraph}
e da representação por matriz de adjacência. \\
\hline
Otávio Santos & Implementação da representação por lista de adjacência
e da hierarquia de exceções. \\
\hline
Arthur Luis & Programa de demonstração, testes e
redação do relatório. \\
\hline
\end{tabular}
\end{table}

\end{document}
